\documentclass[11pt]{article}
\usepackage[margin=1in]{geometry}
\usepackage{amsmath,amssymb,amsthm}
\title{Disproof of Conjecture 00000005350}
\author{earthking11}
\date{October 6, 2026}
\newtheorem{proposition}{Proposition}
\begin{document}
\maketitle
\section*{Markov moves preserve closure}
There are two elementary Markov moves on braids.  Conjugation replaces
$\beta$ by $\alpha\beta\alpha^{-1}$, and stabilization or destabilization
replaces $\beta\in B_n$ by $\beta\sigma_n^{\pm1}\in B_{n+1}$ or conversely.
In both cases the braid closures are isotopic.  This is the preservation
direction of Markov's theorem.

Suppose that a braid $\beta$ is taken to $\gamma$ by one Markov move and
$\gamma$ to $\delta$ by a second.  Closure preservation gives
\[
 \widehat\beta\simeq\widehat\gamma
 \quad\text{and}\quad
 \widehat\gamma\simeq\widehat\delta.
\]
Transitivity of isotopy yields $\widehat\beta\simeq\widehat\delta$.

\begin{proposition}
An explicit pair of Markov moves cannot realize two braids with different
closed-knot types.
\end{proposition}
\begin{proof}
The two applications of closure preservation and transitivity above prove the
closures isotopic, contradicting the requirement that they differ.
\end{proof}

This argument does not need to inspect the categorified invariants.  Whatever
those invariants do, Markov moves themselves cannot be the advertised
mechanism for separating closures.  In fact the same proof applies to every
finite sequence of Markov moves.

The Lean project formalizes the preservation-and-transitivity argument for an
arbitrary braid type, closure type, and Markov-move relation.
\end{document}
